% Macros used only by the printed (xelatex) version.
%
% It starts with dummy definitions that make the pdf compiler ignore the
% macros provided by the plasTeX blueprint plugin, which only make sense for
% the web version.

% Dummy macros that make sense only for web version.
\newcommand{\lean}[1]{}
\newcommand{\discussion}[1]{}
\newcommand{\leanok}{}
\newcommand{\mathlibok}{}
\newcommand{\notready}{}
\newcommand{\alsoIn}[1]{}
% Make sure that arguments of \uses and \proves are real labels, by using invisible refs:
% latex prints a warning if the label is not defined, but nothing is shown in the pdf file.
% It uses LaTeX3 programming, this is why we use the expl3 package.
\ExplSyntaxOn
\NewDocumentCommand{\uses}{m}
 {\clist_map_inline:nn{#1}{\vphantom{\ref{##1}}}%
  \ignorespaces}
\NewDocumentCommand{\proves}{m}
 {\clist_map_inline:nn{#1}{\vphantom{\ref{##1}}}%
  \ignorespaces}
\ExplSyntaxOff

% ---------------------------------------------------------------------------
% Source pointers are long monospace words (e.g.
% Lean4Lean/Verify/Typing/Lemmas.lean:123) and xelatex has no way to break
% them, which produces overfull boxes whenever a node title plus an
% attribution badge plus a \srcloc share the first line of a theorem body.
% Allow a line break after each path separator.  The web version overrides
% \srcloc separately (macros/web.tex), so this stays print-only.
% ---------------------------------------------------------------------------
\ExplSyntaxOn
\tl_new:N \l__blueprint_srcloc_tl
\RenewDocumentCommand{\srcloc}{mm}
  {
    \tl_set:Nn \l__blueprint_srcloc_tl {#1}
    \tl_replace_all:Nnn \l__blueprint_srcloc_tl { / } { / \allowbreak }
    \texttt{ \tl_use:N \l__blueprint_srcloc_tl : #2 }
  }
\ExplSyntaxOff
